\documentclass[12pt,a4paper]{article}
\usepackage[utf8]{inputenc}
\usepackage[T1]{fontenc}
\usepackage[french]{babel}
\usepackage{amsmath, amssymb}
\usepackage{geometry}
\usepackage{array}
\usepackage{longtable}
\usepackage{enumitem}
\geometry{margin=2.5cm}

\title{Formules de dimensionnement SolarEase}
\author{}
\date{}

\begin{document}
\maketitle

\section*{Introduction}
Ce document regroupe les principales formules utilisées dans le service de dimensionnement SolarEase, ainsi que les paramètres à valider avec la société afin d'améliorer la précision des projets.

\section{Surface utile}
\[
\text{usableArea} = \text{roofArea} \times 0.7
\]
Le coefficient $0.7$ représente la surface réellement exploitable.

\section{Nombre de panneaux}
\[
\text{panelCount} = \left\lfloor \frac{\text{usableArea}}{\text{panelAreaM2}} \right\rfloor
\]
Cette formule dépend de la surface d'un panneau et de la surface utile disponible.

\section{Puissance totale installée}
\[
\text{totalCapacityKw} = \frac{\text{panelCount} \times \text{panelPowerW}}{1000}
\]

\section{Production annuelle}
\subsection{Avec PVGIS}
\[
\text{estimatedProduction} = \text{PVGIS}(lat, lon, peakpower, angle, aspect)
\]

\subsection{Sans PVGIS}
\[
\text{estimatedProduction} = \text{totalCapacityKw} \times 1600 \times \text{orientationFactor} \times \text{inclinationFactor}
\]

\section{Facteur d'orientation}
\begin{itemize}[leftmargin=1.5cm]
    \item Sud = $1.0$
    \item Sud-Est / Sud-Ouest = $0.95$
    \item Est / Ouest = $0.85$
    \item Nord-Est / Nord-Ouest = $0.70$
    \item Nord = $0.60$
\end{itemize}

\section{Facteur d'inclinaison}
\[
\text{inclinationFactor} = 1 - \left|30 - \text{inclination}\right| \times 0.005
\]
avec :
\[
\text{inclinationFactor} \ge 0.5
\]

\section{Conversion orientation métier vers PVGIS}
\begin{itemize}[leftmargin=1.5cm]
    \item South = $0^\circ$
    \item South-East = $-45^\circ$
    \item East = $-90^\circ$
    \item North-East = $-135^\circ$
    \item North = $180^\circ$
    \item North-West = $135^\circ$
    \item West = $90^\circ$
    \item South-West = $45^\circ$
\end{itemize}

\section{Coût matériel}
\[
\text{materialCost} = (\text{panelCount} \times \text{panelPrice}) + \text{inverterPrice}
\]

\section{Main-d'œuvre}
\[
\text{installationLabor} = \text{materialCost} \times 0.15
\]

\section{Coût total du projet}
\[
\text{totalProjectCost} = \text{materialCost} + \text{installationLabor}
\]

\section{Économies annuelles}
\[
\text{annualSavings} = \text{estimatedProduction} \times 0.28
\]

\section{Économies mensuelles}
\[
\text{monthlySavings} = \frac{\text{annualSavings}}{12}
\]

\section{CO2 évité}
\[
\text{co2Savings} = \text{estimatedProduction} \times 0.6
\]

\section{Investissement initial financier}
\[
\text{initialInvestment} =
\begin{cases}
\text{estimatedCost} & \text{si disponible} \\
\text{totalCapacityKw} \times 2500 & \text{sinon}
\end{cases}
\]

\section{Dégradation annuelle des panneaux}
\[
\text{currentProduction}_{n+1} = \text{currentProduction}_n \times (1 - 0.005)
\]

\section{Hausse annuelle du prix de l'électricité}
\[
\text{currentElectricityPrice}_{n+1} = \text{currentElectricityPrice}_n \times (1 + 0.04)
\]

\section{Cash flow cumulé}
\[
\text{cashFlow}_0 = -\text{initialInvestment}
\]
\[
\text{cashFlow}_n = \text{cashFlow}_{n-1} + \text{yearSavings}
\]

\section{Temps de retour sur investissement}
Quand le cash flow devient positif :
\[
\text{payback} = (n-1) + \frac{|\text{previousCashFlow}|}{\text{yearSavings}}
\]

\section{ROI}
\[
\text{ROI} = \frac{\text{totalNetSavings}}{\text{initialInvestment}} \times 100
\]

\section{Formules Night Panel}
\subsection{Stockage total}
\[
\text{totalStorageKwh} = \text{panelCount} \times \text{storagePerPanel}
\]

\subsection{Production journalière}
\[
\text{dailyProductionKwh} = \frac{\text{annualProduction}}{365}
\]

\subsection{Consommation de jour}
\[
\text{daytimeConsumption} = \text{dailyConsumption} \times 0.4
\]

\subsection{Consommation de nuit}
\[
\text{nighttimeConsumption} = \text{dailyConsumption} \times 0.6
\]

\subsection{Surplus de jour}
\[
\text{daytimeSurplus} = \max(0, \text{dailyProductionKwh} - \text{daytimeConsumption})
\]

\subsection{Énergie stockée réellement}
\[
\text{actualStored} = \min(\text{daytimeSurplus}, \text{totalStorageKwh})
\]

\subsection{Couverture nocturne}
\[
\text{nightCoverage} = \min(\text{actualStored}, \text{nighttimeConsumption})
\]

\subsection{Taux d'autoconsommation}
\[
\text{selfConsumptionRate} = \min\left(1, \frac{\text{totalSelfConsumed}}{\text{dailyConsumption}}\right)
\]

\subsection{Taux de couverture nocturne}
\[
\text{nightCoverageRate} = \min\left(1, \frac{\text{nightCoverage}}{\text{nighttimeConsumption}}\right)
\]

\subsection{Économies annuelles améliorées}
\[
\text{enhancedAnnualSavings} = \text{totalSelfConsumed} \times 365 \times 0.280
\]

\subsection{Économies liées au surplus réseau}
\[
\text{gridSavings} = \text{gridSurplus} \times 365 \times 0.10
\]

\subsection{Économies annuelles totales}
\[
\text{totalAnnualSavings} = \text{enhancedAnnualSavings} + \text{gridSavings}
\]

\section{Paramètres à valider avec la société}
\begin{itemize}[leftmargin=1.5cm]
    \item Coefficient de surface utile
    \item Productivité moyenne kWh/kWp/an
    \item Facteurs d'orientation
    \item Facteurs d'inclinaison
    \item Prix moyen par kW installé
    \item Prix du kWh économisé
    \item Taux de dégradation annuel
    \item Taux d'inflation de l'électricité
    \item Taux de main-d'œuvre
    \item Facteur CO2
\end{itemize}

\end{document}
